% ============================================================
% 摘要
% ============================================================
\begin{abstract}
深度伪造取证除了回答“真假”，还要回答\emph{“这张伪造图沿用了哪一张真实图的背景”}，
并且往往需要精确到帧——只有定位到具体时刻，才能回到原始素材核对。
本文以“背景即来源指纹”为出发点处理该任务，并让方法由任务性质推出，而不是堆叠模块。

我们先把任务性质翻译成设计约束：背景是唯一稳定的信号；判别信息集中在低频空间布局；
同一视频内的差异微小而局部；几何退化会破坏像素级对应；库中可能没有来源。
由此得到六条约束——背景优先、低频友好、局部成对比较、几何显式校验、按条件启用、
检索与拒识共用同一相似度——并逐条落成 \aura{}：全局编码器初筛、双路局部网格描述子与
可微几何重排、带门控与几何回退的两阶段打分，以及基于 top-$1$ 相似度的开集拒识。
同时固定图库$/$查询构造、帧级$/$视频级判据、只作用于查询的确定性退化（光度六档 $+$ 几何六档），
并补齐“库中可能无源”的闭集$/$开集协议。

按此设计得到的读数与约束逐条对应。标准协议下，干净档帧级由初筛的 $60.00\%$ 升至 $67.67\%$，
几何七档均值 $66.90\%$，与经典 SIFT$+$RANSAC 的 $67.57\%$ 相当；光度退化几乎不改变任何读数，
几何退化才是真正的评测条件。为回答“帧级是否只是从同视频抽出的 $8$ 帧里挑一张”这一质疑，
我们进一步构造稠密帧库（$30$ 条视频各连续 $100$ 帧、逐帧检测人脸），把命中定义为\emph{帧相似度区间}：
\code{clean} 档三种方法在 $p90$ 区间下为 $98.9\%/94.1\%/90.4\%$（同视频随机基线 $23.2\%$），
而几何重排的增益只出现在几何退化档（\code{crop0.8} 的 $T1_{\text{loc}}$ 由 $5.6\%$ 升至 $19.4\%$，
\code{crop0.7off} 由 $3.5\%$ 升至 $12.4\%$，视频级由 $62.2\%$ 升至 $74.1\%$）；
在干净与光度档，全局外观已足够。这一对照限定了旧协议结论的适用范围，
也说明几何校验应当按条件启用而非无条件叠加。

跨数据集证据（Celeb-DF v2 真实换脸，$589$ 库 / $680$ 查询：干净档 $98.82\%$、五档均值 $93.06\%$）
与跨伪造类型的拼接压力测试表明，“以背景为来源指纹”的建模在真实素材上同样成立。
上述协议、因果链与读数一并公开，可作为后续溯源研究的对照起点。

\vspace{2mm}
\noindent\textbf{关键词：}深度伪造；主动检索；来源溯源；无攻击监督学习；退化鲁棒基准；开集拒识；可微几何验证
\end{abstract}

\newpage
